\documentclass[10pt,a4paper]{amsart}
\usepackage[margin=19mm]{geometry}
\usepackage{amsmath,amssymb,mathtools}
\usepackage[scr=rsfs,cal=euler]{mathalfa}
\usepackage{xfrac}
\usepackage[unicode,hidelinks,bookmarks=false]{hyperref}
\newcommand{\ie}{i.e.\ }
\newcommand{\cf}{cf.\ }
\newcommand{\resp}{respectively\ }
\newcommand{\Subsection}{\subsection}
\providecommand{\nobreakdash}{\nobreak\hbox{-}}
\setlength{\emergencystretch}{2em}
\multlinegap0pt
\usepackage{amssymb}
\usepackage{xcolor}
\usepackage{float}
\usepackage{mathtools}
\usepackage{stackengine}
\usepackage[shortlabels]{enumitem}
\newtheorem{prop}{Proposition}
\newtheorem{lem}{Lemma}
\DeclareMathOperator*{\supess}{ess\,sup}
\title{Degenerate parabolic equations in divergence form: fundamental solution and Gaussian bounds}
\author{Khalid Baadi}
\date{Source: arXiv:2503.07569v3 (CC BY 4.0)}
\begin{document}
\maketitle

\noindent\emph{Source and licence.} Degenerate parabolic equations in divergence form: fundamental solution and Gaussian bounds. Version arXiv:2503.07569v3 (CC BY 4.0), distributed by arXiv under the Creative Commons Attribution 4.0 International licence, http://creativecommons.org/licenses/by/4.0/, as recorded in the arXiv licence metadata for this exact version and shown on the abstract page https://arxiv.org/abs/2503.07569v3. Full licence notice: \texttt{licenses/arxiv-2503.07569-CC-BY-4.0.txt}.\par\smallskip

\noindent\textit{Editorial scope.} Counted unit: the article's prop ThmConverseHK (source0.tex lines 827--829). The article's complete original proof of it is delivered with it (source0.tex lines 830--935, one \texttt{proof} environment of 106 lines). No mathematics is written, repaired or re-derived: the statement and the proof are the article's own lines, in the article's own notation, and no retained line is substituted or re-flowed.  Retained uncounted as the source-local context the proof needs: lines 242--244; lines 377--400; lines 482--484; lines 489--495; lines 546--548; lines 598--600; lines 603--609; lines 617--626.  Nothing was omitted from any retained span, so the package carries no omission record.  No retained line contains a citation, so the wrapper carries no bibliography and no bibliography entry.  No retained line contains a figure environment, an image-inclusion command or drawing code, so no image is delivered and the package declares no asset dependency.  Editorial formatting only, in addition: an AMS \texttt{amsart} preamble that restates the article's own declarations and macros used by the retained lines, namely the environments \texttt{prop}, \texttt{lem} on separate counters, together with the standard packages those lines need.  The retained environments print in this document as Proposition 1, Lemma 1 in order of appearance, whereas the article prints the same items with the numbering of its own full document.  The complete native source of the article as distributed in the arXiv e-print for this exact version is bundled unchanged as \texttt{sources/174519/source0.tex}.
\begin{equation}\label{DoublingMuck}
    \omega(2Q)\leq D \omega(Q),
\end{equation}

\begin{prop}[Cauchy problem on $(0,\mathfrak{T})$]\label{prop: Pb de Cauchy borné} Let $ f \in L^2((0,\mathfrak{T});L^2_\omega(\mathbb{R}^n)^n)$, $g\in L^{\rho'}((0,\mathfrak{T});L^2_\omega(\mathbb{R}^n))$ and $\psi \in L^2_\omega(\mathbb{R}^n)$. Then there exists a unique $u \in L^1((0,\mathfrak{T});H^1_\omega(\mathbb{R}^n))$ with $\int_0^\mathfrak{T}\|\nabla_x u(t)\|^2_{2,\omega}\, \mathrm{d}t<\infty$ if $\mathfrak{T}<\infty$ and $u \in L^1_{\mathrm{loc}}((0,\infty);H^1_\omega(\mathbb{R}^n))$ with $\int_0^\infty\|\nabla_x u(t)\|^2_{2,\omega}\, \mathrm{d}t<\infty$ if $\mathfrak{T}=\infty$ solution to the Cauchy problem 
\begin{align*}
\left\{
    \begin{array}{ll}
        \mathcal{H}u  =  - \omega^{-1} \mathrm{div}_x(\omega f) +(-\Delta_\omega)^{\beta/2}g \ \mathrm{in} \  \mathcal{D}'((0,\mathfrak{T})\times \mathbb{R}^n), \\
        u(t) \rightarrow \psi \  \mathrm{ in } \ \mathcal{D'}(\mathbb{R}^n) \ \mathrm{as} \ t \rightarrow 0^+.
    \end{array}
\right.
\end{align*} 
Moreover, $u \in C([0,\mathfrak{T}];L^2_\omega(\mathbb{R}^n))$ with $u(0)=\psi$, $\lim_{t \to \infty} u(t)=0$ if $\mathfrak{T}=\infty$ (by convention, set $u(\infty)=0$), $t \mapsto \| u(t)  \|^2_{2,\omega}$ is absolutely continuous on $[0,\mathfrak{T}]$ and we can write the energy equalities. Furthermore, we have $(-\Delta_\omega)^{\alpha/2} u \in  L^r((0,\mathfrak{T});L^2_\omega(\mathbb{R}^n))$ for any $\alpha \in (0,1]$ with $r=\frac{2}{\alpha} \in [2,\infty)$ with
\begin{align*}
            \sup_{t\in [0,\mathfrak{T}]} \| u(t) \|_{2,\omega}+ \| (-\Delta_\omega)^{\alpha /2} u\|_{L^r((0,\mathfrak{T});L^2_\omega(\mathbb{R}^n))} 
            \leq C  ( \| f  \|_{L^2((0,\mathfrak{T});L^2_\omega(\mathbb{R}^n)^n)}+ \left \| g \right \|_{L^{\rho'}((0,\mathfrak{T});L^2_\omega(\mathbb{R}^n))}+  \| \psi  \|_{2,\omega}  ),
\end{align*} 
where $C=C(M,\nu,\rho)>0$ is a constant. Lastly,  for all $t \in [0,\mathfrak{T}]$, we have the following representation of $u$ (by convention, set $\Gamma(\infty,s)=0$ if $\mathfrak{T}=\infty$): 
\begin{align}\label{rep}
    u(t) = \Gamma(t,0)a  + \int_{0}^{t} \Gamma(t,\tau) (- \omega^{-1} \mathrm{div}_x(\omega f))(\tau)\ \mathrm d \tau  + \int_{0}^{t} \Gamma(t,\tau) (-\Delta_\omega)^{\beta/2}g(\tau)\ \mathrm d \tau,
\end{align}
where the two integrals are weakly defined in $L^2_\omega(\mathbb{R}^n)$. More precisely,  for all $\Tilde{\psi}\in L^2_\omega(\mathbb{R}^n)$, we have the equality with absolutely converging integrals
    \begin{align*}
        \langle u(t), \Tilde{\psi} \rangle_{2,\omega} = \langle \Gamma(t,0)\psi, \Tilde{\psi} \rangle_{2,\omega}+ \int_{0}^{t} \langle f(\tau), \nabla_x \Tilde{\Gamma}(\tau,t)\Tilde{\psi} \rangle_{2,\omega}\  \mathrm d \tau + \int_{0}^{t} \langle  g(\tau),(-\Delta_\omega)^{\beta/2}\Tilde{\Gamma}(\tau,t) \Tilde{\psi} \rangle_{2,\omega} \ \mathrm d \tau.
    \end{align*}
When $\rho=\infty$, or equivalently $\beta = 0$, then the last integral in \eqref{rep} converges also strongly in $L^2_\omega(\mathbb{R}^n)$, \textit{i.e.} in the Bochner sense.
\end{prop}

\begin{equation}\label{sup pour sol faibles}
        \sup_{t \in I} \left ( \supess_{\mathcal{O'}} \left | u(t,\cdot) \right | \right ) = \supess_{I \times \mathcal{O'}} \left | u \right |,
\end{equation}

\begin{lem}\label{lem:Caccioppoli}
    Let $(t_0,x_0)\in \mathbb{R}\times \mathbb{R}^n$ and $R>0$. Let $u$ be a local weak solution of $\mathcal{H}u=0$ on a neighborhood (of type $I \times \mathcal{O}$) of $Q_{2R}(t_0,x_0)$. Then, we have the following localized energy inequality:
    \begin{equation*}
       \sup_{t_0-R'^2 \leq t \leq t_0 } \left\|u(t) \right\|^2_{L^2_\omega(B(x_0,R'))} + \int_{Q_{R'}(t_0,x_0)} \left | \nabla_x u \right |^2 \ \mathrm d\mu \leq \frac{C}{(2R-R')^2} \int_{Q_{2R}(t_0,x_0)} \left |  u \right |^2 \ \mathrm d\mu,
    \end{equation*}
    for all $R'\in (0,2R)$, where $C=C(n,M,\nu)>0$ is a constant.
\end{lem}

\begin{equation}\label{Variante Caccioppoli}
    \int_{\Sigma} | \nabla_x (\Tilde{\Gamma}(\cdot,t_0)\psi) |^2 \ \mathrm{d\mu} \leq \frac{C}{\min(\rho_2 - \rho_1, \rho_4 - \rho_3)} \frac{1}{R^2} \int_{\Sigma'} | \Tilde{\Gamma}(\cdot,t_0)\psi |^2 \ \mathrm{d\mu},
\end{equation}

\begin{equation}\label{pgb}
    \left | \Gamma(t,x;s,y) \right | \leq \frac{K_0}{\sqrt{\omega_{t-s}(x)}\sqrt{\omega_{t-s}(y)}} e^{-k_0 \frac{|x-y|^2}{t-s}},
\end{equation}

\begin{lem}\label{lem:Cruz-Uribe et Rios}
    The factor $\frac{1}{\sqrt{\omega_{t-s}(x)}\sqrt{\omega_{t-s}(y)}}$ appearing in \eqref{pgb} above may be replaced by one of 
    \begin{equation*}
        \frac{1}{\omega_{t-s}(x) }, \ \ \frac{1}{\omega_{t-s}(y) }, \ \ \frac{1}{\max(\omega_{t-s}(x) ,\omega_{t-s}(y))},
    \end{equation*}
    and the constants $K_0$ and $k_0$ in \eqref{pgb} are replaced respectively by $\Tilde{K}_0=\Tilde{K}_0(K_0,k_0,D)>0$ and $\frac{k_0}{2}$.
\end{lem}

\begin{prop}[Properties of the generalized fundamental solution]\label{prop:solfundgeneralisée}
    $\mathcal{H}$ has Gaussian upper bounds if and only if $\mathcal{H}^\star$ does. In this case, the following properties hold for all $t>s$:
    \begin{enumerate}
        \item (Adjointess property) For almost every $(x,y) \in \mathbb{R}^{2n}$, we have
$$ \Tilde{\Gamma}(s,y;t,x) = \overline{\Gamma(t,x;s,y)}. $$
        \item (Chapman-Kolmogorov identities) If $t>r>s$, then for almost every $(x,y)\in \mathbb{R}^{2n}$, we have
        $$ \int_{\mathbb{R}^n} \Gamma(t,x;r,z) \Gamma(r,z;s,y) \ \mathrm{d}\omega(z)= \Gamma(t,x;s,y). $$
        \item (Conservation property) For almost every $x\in \mathbb{R}^n$, we have $$ \int_{\mathbb{R}^n} \Gamma(t,x;s,y) \ \mathrm{d}\omega(y)=1. $$
    \end{enumerate}
\end{prop}

\begin{prop}\label{ThmConverseHK}
    If $\mathcal{H}$ has Gaussian upper bounds, then $\mathcal{H}$ and $\mathcal{H}^\star$ satisfy Moser’s $L^2$-$L^\infty$ estimates.
\end{prop}

\begin{proof}
The proof is the same for both $\mathcal{H}$ and $\mathcal{H}^\star$ and we only prove the result for $\mathcal{H}$. We fix $R>0$, $(t_0,x_0) \in \mathbb{R}^{1+n}$, and let $u$ be a local weak solutions of $\mathcal{H}u=0$ on a neighborhood of $Q_{2R}(t_0,x_0)$.
Using \eqref{sup pour sol faibles}, we need to prove that 
\begin{equation}\label{JJ}
      \sup_{t \in (t_0-R^2, t_0]} \left ( \supess_{B(x_0,R)} \left | u(t,\cdot) \right | \right )   \leq B \left ( \frac{1}{\mu(Q_{2R}(t_0,x_0))} \int_{Q_{2R}(t_0,x_0)} \left | u \right |^2 \ \mathrm d\mu\right  )^{1/2},
\end{equation}
where $B$ is a constant depending only on $K_0,k_0,\nu,M$, $D$ and $n$. To do so, let $\zeta$ be a non negative smooth function such that $\zeta = 1$ on $Q_{\frac{3}{2}R}(t_0,x_0)$, $\zeta = 0$ outside $Q_{\frac{7}{4}R}(t_0,x_0)$ and verifying $$\left \| \partial_t \zeta  \right \|_{L^\infty(\mathbb{R}^{1+n})}+\left \| \nabla_x \zeta  \right \|^2_{L^\infty(\mathbb{R}^{1+n})} \leq \frac{c(n)}{R^2}.$$
Then $v := \zeta u$ is the unique solution $v \in L^2((s,t);H^1_\omega(\mathbb{R}^n))$ to the  following Cauchy problem
\begin{align*}
\left\{
    \begin{array}{ll}
        \mathcal{H} v = (\partial_t \zeta) u - \omega^{-1}A(t,\cdot)\nabla_x u \cdot \nabla_x \zeta - \omega^{-1}\mathrm{div}_x(A(t,\cdot) u \nabla_x \zeta) \ \mathrm{in} \  \mathcal{D}'((t_0-4R^2,t_0) \times \mathbb{R}^n), \\
        v(\tau) \rightarrow 0 \  \mathrm{ in } \ \mathcal{D'}(\mathbb{R}^n) \ \mathrm{as} \ \tau \rightarrow (t_0-4R^2)^+
    \end{array}
\right.
\end{align*} 
By uniqueness and linearity, we write $v=v_1+v_2+v_3$ where $v_k$ is the solution to the above Cauchy problem considering only the $k^{th}$ term in the right-hand side of the first equation. We fix $ t\in (t_0-R^2,t_0]$. We have $v(t)=v_1(t)+v_2(t)+v_3(t)$ in $L^2_\omega(\mathbb{R}^n)$. Hence, for almost every $x \in B(x_0,R)$,
$$u(t,x)=v(t,x)=v_1(t,x)+v_2(t,x)+v_3(t,x).$$
\underline{\textit{Estimate of $\left\| v_1(t,\cdot) \right\|_{L^\infty(B(x_0,R))}$}:} since $(\partial_t \zeta) u$ is compactly supported in time and space, the representation in Proposition \ref{prop: Pb de Cauchy borné} can be written pointwisely as follows : for almost every $x \in B(x_0,R)$, we have
\begin{equation*}
    v_1(t,x)= \int_{t_0-4R^2}^{t}\left ( \Gamma(t,s)(\partial_t \zeta u)(s) \right )(x) \ \mathrm{d}s = \int_{t_0-4R^2}^{t} \int_{\mathbb{R}^n} \Gamma(t,x;s,y)(\partial_t \zeta u)(s,y) \ \mathrm{d} \omega(y) \ \mathrm ds.
\end{equation*}
Since $\partial_t \zeta = 0 $ on $Q_{\frac{3}{2}R}(t_0,x_0)$ and outside $Q_{\frac{7}{4}R}(t_0,x_0)$, we have by Cauchy-Schwarz inequality, 
\begin{equation*}
    \left | v_1(t,x) \right |\leq \left ( \int_{\Sigma} \left | \Gamma(t,x;s,y) \right |^2 \ \mathrm d\mu(s,y) \right )^{1/2}  \left ( \int_{Q_{\frac{7}{4}R}(t_0,x_0) } \left | (\partial_t \zeta u)(s,y) \right |^2 \ \mathrm d\mu(s,y) \right )^{1/2} ,
\end{equation*}
where $\Sigma:= (s<t)\cap(Q_{\frac{7}{4}R}(t_0,x_0) \setminus Q_{\frac{3}{2}R}(t_0,x_0))$. Therefore, for almost every $x \in B(x_0,R)$,
\begin{equation}\label{1 converse HK}
    \left | v_1(t,x) \right |\leq \frac{c(n)}{R^2}\left ( \int_{\Sigma} \left | \Gamma(t,x;s,y) \right |^2 \ \mathrm d\mu(s,y) \right )^{1/2}  \left ( \int_{Q_{2R}(t_0,x_0) } \left | u \right |^2 \ \mathrm d\mu \right )^{1/2}.
\end{equation}
\underline{\textit{Estimate of $\left\| v_2(t,\cdot) \right\|_{L^\infty(B(x_0,R))}$}:} likewise, since $- \omega^{-1}A(t,\cdot)\nabla_x u \cdot \nabla_x \zeta$ is compactly supported in time and space, we have by Cauchy-Schwarz inequality and for almost every $x \in B(x_0,R)$,
\begin{align*}
    \left | v_2(t,x) \right |&\leq \left ( \int_{\Sigma} \left | \Gamma(t,x;s,y) \right |^2 \ \mathrm d\mu(s,y) \right )^{1/2}  \left ( \int_{Q_{\frac{7}{4}R}(t_0,x_0) } \left | \left ( - \omega^{-1}A(t,\cdot)\nabla_x u \cdot \nabla_x \zeta \right )(s,y) \right |^2 \ \mathrm d\mu(s,y) \right )^{1/2}\\
    &\leq \frac{M \times c(n)}{R} \left ( \int_{\Sigma} \left | \Gamma(t,x;s,y) \right |^2 \ \mathrm d\mu(s,y) \right )^{1/2}  \left ( \int_{Q_{\frac{7}{4}R}(t_0,x_0) } \left | \nabla_x u (s,y) \right |^2 \ \mathrm d\mu(s,y) \right )^{1/2}.
\end{align*}
Using Caccioppoli inequality from Lemma \ref{lem:Caccioppoli}, we deduce that for almost every $x \in B(x_0,R)$,
\begin{equation}\label{2 Convers HK}
    \left | v_2(t,x) \right |\leq \frac{C(n,M,\nu)}{R^2} \left ( \int_{\Sigma} \left | \Gamma(t,x;s,y) \right |^2 \ \mathrm d\mu(s,y) \right )^{1/2}  \left ( \int_{Q_{2R}(t_0,x_0) } \left |  u  \right |^2 \ \mathrm d\mu \right )^{1/2}.
\end{equation}
\underline{\textit{Estimate of $\left\| v_3(t,\cdot) \right\|_{L^\infty(B(x_0,R))}$}:} to obtain a similar estimate on $v_3$, we use again the representation in Proposition \ref{prop: Pb de Cauchy borné} to write 
\begin{equation*}
    \langle v_3(t), \phi \rangle_{2,\omega} = \int_{t_0-4R^2}^{t} \langle \omega^{-1} A(s,\cdot)(\nabla_x \zeta)u, \nabla_x \Tilde{\Gamma}(s,t)\phi \rangle_{2,\omega} \ \mathrm{d} s, \ \ \forall \phi \in \mathcal{D}(\mathbb{R}^n).
\end{equation*}
We fix $\phi \in \mathcal{D}(\mathbb{R}^n)$ such that $\mathrm{supp}(\phi)\subset B(x_0,R)$. We write
\begin{align*}
   \langle v_3(t), \phi \rangle_{2,\omega} &= \int_{t_0-4R^2}^{t} \int_{\mathbb{R}^n} \omega^{-1}(y) u(s,y)\hspace{0.05cm} A(s,y) \hspace{0.05cm} \nabla_x \zeta (s,y)   \cdot \overline{\nabla_x \Tilde{\Gamma}(s,t)\phi (y)} \ \mathrm{d} \omega(y) \ \mathrm{d}s\\&= \int_{\Sigma} \omega^{-1} u \hspace{0.05cm}A \hspace{0.05cm} \nabla_x \zeta \cdot \overline{\nabla_x \Tilde{\Gamma}(\cdot,t)\phi} \ \mathrm{d} \mu.
\end{align*}
Hence,
\begin{equation*}
    \left |  \langle v_3(t), \phi \rangle_{2,\omega} \right |\leq \frac{M\times c(n)}{R}  \left ( \int_{Q_{\frac{7}{4}R}(t_0,x_0) } \left |  u  \right |^2 \ \mathrm d\mu \right )^{1/2}  \left ( \int_{\Sigma}  | \nabla_x \Tilde{\Gamma}(s,t)\phi (y) |^2 \ \mathrm d\mu(s,y)\right )^{1/2}.
\end{equation*}
We define $\Sigma':=Q_{4R}(t,x_0)\setminus Q_{\sqrt{\frac{9}{8}}R}(t,x_0)$. We have $\Sigma \subset Q_{3R}(t,x_0)\setminus Q_{\sqrt{\frac{5}{4}}R}(t,x_0) \subset \Sigma'$. Using the variant of the Caccioppoli inequality \eqref{Variante Caccioppoli}, we have 
\begin{equation*}
    \int_{\Sigma}  | \nabla_x \Tilde{\Gamma}(s,t)\phi (y)   |^2 \ \mathrm d\mu(s,y) \leq \frac{C(n,M,\nu)}{R^2} \int_{\Sigma'}   |  \Tilde{\Gamma}(s,t)\phi (y)   |^2 \ \mathrm d\mu(s,y),
\end{equation*}
Hence, there is a constant $\Tilde{C}=\Tilde{C}(n,M,\nu)>0$ such that
\begin{align}\label{QQQQQ}
    \left |  \langle v_3(t), \phi \rangle_{2,\omega} \right |\leq \frac{\Tilde{C}}{R^2}  \left ( \int_{Q_{2R}(t_0,x_0) } \left |  u  \right |^2 \ \mathrm d\mu \right )^{1/2}  \left ( \int_{\Sigma'}  |  \Tilde{\Gamma}(s,t)\phi (y)  |^2 \ \mathrm d\mu(s,y)\right )^{1/2}.
\end{align}
Since $\Gamma(t,s)$ represented by the kernel $\Gamma(t,x;s,y)$, its adjoint $\Tilde{\Gamma}(s,t)$ is also represented by the kernel $\Tilde{\Gamma}(s,y;t,x)= \overline{\Gamma(t,x;s,y)}$ by Proposition \ref{prop:solfundgeneralisée}. We can therefore write 
\begin{align*}
    \int_{\Sigma'} |  \Tilde{\Gamma}(s,t)\phi (y)    |^2 \ \mathrm d\mu(s,y) &= \int_{\Sigma'} \left | \int_{B(x_0,R)}  \overline{\Gamma}(t,z;s,y) \phi(z)\ \mathrm d\omega(z) \right |^2 d \mu(s,y)
    \\ & \leq \supess_{z \in B(x_0,R)} \left ( \int_{\Sigma'}  | \Gamma(t,z;s,y)  |^2 \ \mathrm d\mu(s,y)  \right ) \left \| \phi \right \|^2_{L^1_\omega(B(x_0,R))}.
\end{align*}
Using this in \eqref{QQQQQ}, we have 
\begin{equation*}
    \left |  \langle v_3(t), \phi \rangle_{2,\omega} \right |\leq \frac{\Tilde{C}}{R^2}  \left ( \int_{Q_{2R}(t_0,x_0) } \left |  u  \right |^2 \ \mathrm d\mu \right )^{1/2} \left (\supess_{z \in B(x_0,R)} \int_{\Sigma'}  | \Gamma(t,z;s,y)     |^2 \ \mathrm d\mu(s,y)  \right )^{1/2} \left \| \phi \right \|_{L^1_\omega(B(x_0,R))}.
\end{equation*}
The inequality above is valid for all $L^2$-functions on $B(x_0,R)$. Using the Lebesgue differentiation theorem, we deduce that for almost every $x \in B(x_0,R)$,
\begin{equation}\label{3 Converse HK}
    \left |  v_3(t,x) \right | \leq  \frac{\Tilde{C}}{R^2} \left (\supess_{z \in B(x_0,R)} \int_{\Sigma'} | \Gamma(t,z;s,y)     |^2 \ \mathrm d\mu(s,y)  \right )^{1/2} \left ( \int_{Q_{2R}(t_0,x_0) } \left |  u  \right |^2 \ \mathrm d\mu \right )^{1/2}  .
\end{equation}
Combining \eqref{1 converse HK}, \eqref{2 Convers HK}, \eqref{3 Converse HK}, the assumption \eqref{pgb} on the kernel and Lemma \ref{lem:Cruz-Uribe et Rios}, we deduce that for almost every $x \in B(x_0,R)$
\begin{equation}\label{LL}
    \left |  u(t,x) \right | \leq  \frac{C(n,M,\nu,K_0,k_0)}{R^2} \left (\supess_{z \in B(x_0,R)} \int_{\Sigma'} \frac{1}{\omega_{t-s}(z)^2} e ^{-k_0 \frac{\left | z-y \right |^2}{t-s}} \ \mathrm d\mu(s,y)  \right )^{1/2} \left ( \int_{Q_{2R}(t_0,x_0) } \left |  u  \right |^2 \ \mathrm d\mu \right )^{1/2},
\end{equation}
since $\Sigma \subset \Sigma'$. We claim that the following estimate holds for all $z \in B(x_0,R)$.
\begin{equation}\label{Poids}
    \int_{\Sigma'} \frac{1}{\omega_{t-s}(z)^2} e ^{-k_0 \frac{\left | z-y \right |^2}{t-s}} \ \mathrm d\mu(s,y) \leq C(k_0,D) \frac{R^2}{\omega(B(x_0,2R))}.
\end{equation}
Then the estimate \eqref{JJ} follows from \eqref{LL}.\\
It remains to prove the claim \eqref{Poids}. Using a change of variable in $s$, we have for all $z \in B(x_0,R)$,
\begin{align*}
    \int_{\Sigma'} \frac{1}{\omega_{t-s}(z)^2} e ^{-k_0 \frac{\left | z-y \right |^2}{t-s}} \ \mathrm d\mu(s,y) &= \int_{0}^{\frac{9}{8}R^2}\int_{\sqrt{\frac{9}{8}}R\leq \left |x_0- y \right |<4R} + \int_{\frac{9}{8}R^2}^{16R^2}\int_{B(x_0,4R)} \frac{\exp(-k_0 \frac{\left |z- y \right |^2}{r})}{\omega_r(z)^2} \omega(y)\ \mathrm dy \mathrm d r \\&:= I +II.
\end{align*}
For $y \in B(x_0,4R)\setminus B(x_0,\sqrt{\frac{9}{8}}R)$, we have $\left | z-y \right | \ge \left | x_0-y \right |-\left | z-x_0 \right | \ge (\sqrt{\frac{9}{8}}-1)R$. Thus, we have
\begin{align*}
    I&= \sum_{k=1}^{+\infty} \int_{ \frac{9}{8}\frac{R^2}{4^{k+1}}}^{ \frac{9}{8}\frac{R^2}{4^k}}\int_{\sqrt{\frac{9}{8}}R\leq \left |x_0- y \right |<4R}  \frac{\exp(-k_0 \frac{\left | z-y \right |^2}{r})}{\omega_r(z)^2} \omega(y)\ \mathrm dy \mathrm d r
    \\&  \leq \sum_{k=1}^{+\infty} \int_{ \frac{9}{8}\frac{R^2}{4^{k+1}}}^{ \frac{9}{8}\frac{R^2}{4^k}}\int_{\sqrt{\frac{9}{8}}R\leq \left |x_0- y \right |<4R}  \frac{e^{-k_0 \left ( \sqrt{\frac{9}{8}}-1 \right )^2\frac{R^2}{r}}}{\omega_r(z)^2} \omega(y)\ \mathrm dy \mathrm d r
    \\& \leq \sum_{k=1}^{+\infty} \int_{ \frac{9}{8}\frac{R^2}{4^{k+1}}}^{ \frac{9}{8}\frac{R^2}{4^k}} \frac{1}{\omega_r(z)^2} \ \mathrm dr \ \omega(B(x_0,4R)) e^{-\frac{8k_0}{9}\left ( \sqrt{\frac{9}{8}}-1 \right )^2 4^k} 
    \\& \leq \sum_{k=1}^{+\infty} \ \frac{1}{\omega(B(z,\sqrt{\frac{9}{8}}\frac{R}{2^{k+1}}))^2} \frac{9}{8}R^2  \left ( \frac{1}{4^k}- \frac{1}{4^{k+1}} \right ) e^{-\frac{8k_0}{9}\left ( \sqrt{\frac{9}{8}}-1 \right )^2 4^k} \omega(B(x_0,4R))
\end{align*}
The doubling property \eqref{DoublingMuck} imply that for all $k\ge0$, $\omega(B(z,\sqrt{\frac{9}{8}}R)) \leq D^{k+1}\omega(B(z,\sqrt{\frac{9}{8}}\frac{R}{2^{k+1}}))$. Thus,
\begin{align*}
    I\leq \left (  \sum_{k=1}^{+\infty} D^{2(k+1)} \left ( \frac{1}{4^k}- \frac{1}{4^{k+1}} \right ) e^{-\frac{8k_0}{9}\left ( \sqrt{\frac{9}{8}}-1 \right )^2 4^k}\right )  \frac{\omega(B(x_0,4R))}{\omega(B(z,\sqrt{\frac{9}{8}}R))^2} \frac{9}{8} R^2.
\end{align*}
As $|x_0-z| \leq R$, we use the doubling property \eqref{DoublingMuck} to deduce that
\begin{equation}\label{I}
     I \leq C(k_0,D) \frac{R^2}{\omega(B(x_0,2R))}.
\end{equation}
Estimating $II$ is straightforward by simply writing  
\begin{equation}\label{II}
    II \leq \frac{\omega(B(x_0,R))}{\omega(B(z,\sqrt{\frac{9}{8}}R))^2} \left ( 16R^2- \frac{9}{8}R^2 \right ) \leq C(D) \frac{R^2}{\omega(B(x_0,2R))},
\end{equation}
by the doubling property \eqref{DoublingMuck} as $|x_0-z| \leq R$. The estimate \eqref{Poids} follows from \eqref{I} and \eqref{II}. 
\end{proof}

\end{document}
